\begin{textsample}{POS Dim 2 – persona\_gemini – Score 22.00 – t195\_gemini.txt}  \label{ex:f2_pos_041}
For Filipino farmers and fishers, \textbf{climate} change has disrupted traditional \textbf{weather} \textbf{knowledge}, threatening their food \textbf{security} and well-being.
They \textbf{face} frequent typhoons, \textbf{floods}, and \textbf{droughts}, with these \textbf{impacts} exacerbated by \textbf{poverty} and poor governance.
The term"\textbf{resilience}"is often used to describe these \textbf{communities}, but this framing can be a critique.
It burdens them with the responsibility to adapt while absolving fossil \textbf{fuel} companies and \textbf{governments} of their own responsibility for the \textbf{crisis}.
In response, a petition was filed in 2015 against major fossil \textbf{fuel} producers, or"carbon majors,"which led to a \textbf{human} \textbf{rights} investigation in 2019.
The \textbf{urgency} of the situation is confirmed by the recent IPCC report, which warns of escalating \textbf{climate} \textbf{risks} and higher mortality rates in vulnerable areas.
To \textbf{limit} warming to 1.5°C, the report stresses the need for \textbf{ecosystem} \textbf{protection}, adaptation measures, addressing \textbf{loss} and damage, and ensuring \textbf{climate} \textbf{justice}.

% matched lemmas: climate, community, crisis, drought, ecosystem, face, flood, fuel, government, human, impact, justice, knowledge, limit, loss, poverty, protection, resilience, right, risk, security, urgency, weather
\end{textsample}
